\section{Soundness, completeness, and a constructive negative result}
\label{sec:theory}

\subsection{Soundness for arbitrary polynomial updates}
\begin{theorem}[Finite closure certificate]\label{thm:closure}
For the machine of Section~\ref{sec:model}, suppose the polynomial identities
\eqref{eq:goalcert}--\eqref{eq:initcert} hold. Then every reachable state
satisfies all its demanded output equalities.
\end{theorem}
\begin{proof}
Induct on the length of an execution. Equation~\eqref{eq:initcert} proves that
all initial generators vanish. Suppose $P_q(x)=0$ and an edge with payload $u$
moves the state to $F_e(x,u)$ at $q'$. Evaluating
\eqref{eq:edgecert} at this particular $x,u$ yields
$P_{q'}(F_e(x,u))=M_e(u)0=0$. Thus the generators vanish after every finite
execution. At its endpoint, \eqref{eq:goalcert} gives $G_q(x)=C_q0=0$.
No assumption that the update is affine, invertible, or terminating beyond the
finite input word was used.
\end{proof}

The quantifier over input words is unbounded. The certificate proves all word
lengths at once; it is not an extrapolation from a bounded execution trace.
The argument also explains why matrices need not be nonnegative: these are
zero equalities, not order inequalities.

\begin{corollary}[Relational fold verification]
Suppose a faithful product-machine translation of two folds has output goals
$o_f(s)-o_h(t)=0$ and admits a closure certificate. Then their outputs agree
on every finite input list in the translated input domain.
\end{corollary}
\begin{proof}
A simultaneous step of the product updates each component on the same payload.
Induction on the input list identifies the product's endpoint with the pair of
fold endpoints. Apply Theorem~\ref{thm:closure} to its output difference.
\end{proof}

\subsection{What the constructed space means}
\begin{lemma}[Path-coefficient characterization]\label{lem:path}
For a control node $q$, the least space $W_q$ is spanned by the coefficient
polynomials of all expressions obtained by pulling a demanded output backward
along a finite path beginning at $q$. Every time step uses its own fresh tuple
of payload indeterminates.
\end{lemma}
\begin{proof}
Length-zero paths contribute precisely the requested goals. Extending a path
backward by one edge first substitutes that edge's update, then extracts
coefficients in its fresh inputs. Coefficient extraction for the other time
steps commutes with this operation because their variables are distinct.
Consequently the span of all path coefficients contains the goals and is
closed under every $T_{e,\alpha}$. Conversely, repeated closure applied to a
goal produces only such coefficients or their rational linear combinations.
Every closed family containing the goals contains them all. This proves both
inclusions and leastness.
\end{proof}

Leastness is a vector-space statement. It does not say this is the smallest
inductive invariant as a set of states, or a basis of the ideal of all reachable
polynomial equalities. The distinction matters for nonlinear updates.

\subsection{Constructing a counterexample instead of merely rejecting a basis}
\begin{lemma}[A finite separating grid]\label{lem:grid}
Let $f\in\Q[u_1,\ldots,u_m]$ be nonzero and let $d_i$ be its degree in $u_i$.
There is a point in $\prod_i\{0,1,\ldots,d_i\}$ at which $f$ is nonzero.
For $m=0$, the grid consists of the unique empty tuple.
\end{lemma}
\begin{proof}
For one variable, a nonzero degree-$d$ polynomial has at most $d$ roots. For
several variables, write $f$ as a polynomial in the last variable. If it
vanished on the whole grid, then at every choice of the earlier coordinates
all its coefficients in the last variable would vanish. Induction on the
number of variables would make every coefficient polynomial zero, contrary to
$f\ne0$.
\end{proof}

Suppose a retained polynomial $p\in P_{q_0}$ has $p(x_0)\ne0$. If its origin
is a requested goal, the empty trace already refutes the property. Otherwise
its origin records an edge $e:q_0\to q'$ and an older polynomial $h\in P_{q'}$
for which $p$ is one payload coefficient of $h(F_e(x,u))$. At the concrete
current state $x_0$, that coefficient is nonzero. The polynomial
$u\mapsto h(F_e(x_0,u))$ is therefore nonzero. Lemma~\ref{lem:grid} supplies a
concrete payload $u_0$ making it nonzero. Execute the edge and repeat with $h$.

Each parent was retained before its child. Following parents must therefore
terminate at a demanded goal. The resulting certificate is a finite sequence
of original edge indices and rational payload tuples, together with the index
of the final output that is nonzero. A separate checker executes the original
updates and verifies the final discrepancy. It does not need the basis,
provenance graph, or root-bound argument.

\begin{proposition}[Negative evidence from a retained origin]
With unrestricted rational payloads and complete original-generator provenance,
a nonvanishing generator at the initial state yields a genuine finite
execution counterexample. A budget cutoff while constructing its payloads
justifies only \unknown{}, not an accepted counterexample.
\end{proposition}
\begin{proof}
At each selected step the preceding paragraph preserves the assertion that
the current retained observation is nonzero. The provenance order is strictly
decreasing. Its last node is an original demanded goal, and all traversed
edges start at the current control node. The constructed trace is therefore
an execution whose endpoint violates that goal. Independent execution replay
checks this final claim directly.
\end{proof}

\subsection{A complete fragment with a finite dimension bound}
\begin{definition}[State-affine update]
An update is state-affine when
\[
 F_e(x,u)=A_e(u)x+b_e(u),
\]
where the entries of $A_e$ and $b_e$ are rational polynomials in the fresh
payloads. It need not be affine in the payloads.
\end{definition}

\begin{theorem}[Exact decision in the state-affine fragment]\label{thm:complete}
Assume all updates are state-affine and every demanded output has state degree
at most $D$. In exact arithmetic, with no resource truncation, observable
closure stabilizes after at most
\[
 B=|Q|\binom{r+D}{D}
\]
independent generator insertions. The procedure either returns a positive
certificate for \eqref{eq:task} or returns a concrete execution counterexample.
\end{theorem}
\begin{proof}
Substitution of state-affine expressions into a degree-$D$ observation does not
increase its degree in $x$. Extracting coefficients in $u$ cannot increase that
degree either. All spaces therefore remain in
$V_D=\{p\in\Q[x]:\deg p\le D\}$, whose monomial basis has cardinality
$N=\binom{r+D}{D}$. Each insertion strictly increases the dimension of one
node's space. There are at most $|Q|N$ such increases, each generating finitely
many coefficient obligations. The finite worklist eventually empties unless
an initial generator is found nonzero.

If an initial generator is nonzero, provenance reconstruction gives a
counterexample. Otherwise all initial generators vanish. At stabilization,
exact linear dependence supplies the matrices in \eqref{eq:goalcert} and
\eqref{eq:edgecert}; hence Theorem~\ref{thm:closure} proves the property.
Equivalently, Lemma~\ref{lem:path} shows that a valid property forces every
path-coefficient generator to vanish initially, since a polynomial vanishing
on every rational payload assignment has zero coefficients.
\end{proof}

\begin{corollary}[A short counterexample word]
If a property in Theorem~\ref{thm:complete}'s fragment is false, the retained
origin construction can produce a counterexample with at most $B-1$ edges.
This bounds the word length, not the bit size of its rational states.
\end{corollary}
\begin{proof}
A nonzero retained initial generator has an ancestry chain ending in an
original goal. Each edge on that chain moves to a strictly earlier retained
generator. At most $B$ generators can be retained in the whole control graph,
so such a chain contains at most $B-1$ coefficient steps.
\end{proof}

The code enforces practical degree and basis bounds. The mathematical theorem
is not a promise that those defaults are large enough for every legal input.
For example, if $B$ exceeds the configured per-node capacity, a valid problem
can still produce \unknown{}.

\subsection{Nonlinear success and nonlinear failure are both informative}
For $s'=s+u$ and $t'=t+3s^2u+3su^2+u^3$, the degree-three observation
$t-s^3$ is conserved. The nonlinear-state fragment succeeds immediately.
But for $x'=x^2$, $x_0=0$, and goal $x=0$, closure generates
\[
 x,\ x^2,\ x^4,\ x^8,\ldots
\]
which are linearly independent as polynomials. The property is true, yet the
implemented degree cap correctly returns \unknown{}.

The ideal generated by $x$ would handle this last example because $x^2=x\cdot x$.
That observation points to a complementary use of Forge's existing
polynomial-multiplier machinery, not an excuse to describe vector-space closure
as a general invariant decision procedure.

\subsection{Cost and specialization}
Let $N=\binom{r+D}{D}$, and suppose every pullback has at most $U$ distinct
payload monomials. Each edge is processed once for every retained generator at
its destination, giving at most $EN$ pullbacks and $ENU$ coefficient
candidates for $E$ edges. A crude dense echelon bound is $O(EUN^3)$ rational
operations, plus the cost of polynomial expansion. Dense space and coordinate
storage requires $O(|Q|N^2)$ coefficients. Coefficient bit growth is not
included in those arithmetic-operation counts.

The advantage of demand direction is that the actual dimensions $k_q$ can be
much smaller than $N$. The delayed example uses two generators out of five
possible affine monomials, and the square-of-sum example uses one out of six
quadratic monomials. These are exact dimension comparisons, not claims of a
measured runtime advantage over a full-lifting implementation.
